%% IEEE 期刊 / 会议通用投稿骨架
%% 依赖 TeX Live 自带 IEEEtran（LPPL）
%% 编译: xelatex 或 pdflatex
\documentclass[conference]{IEEEtran}
% 期刊双栏可改为: \documentclass[journal]{IEEEtran}

\usepackage{cite}
\usepackage{amsmath,amssymb,amsfonts}
\usepackage{algorithmic}
\usepackage{graphicx}
\usepackage{textcomp}
\usepackage{xcolor}

\def\BibTeX{{\rm B\kern-.05em{\sc i\kern-.025em b}\kern-.08em
    T\kern-.1667em\lower.7ex\hbox{E}\kern-.125emX}}

\begin{document}

\title{A Concise Template for IEEE Conferences and Transactions}

\author{
\IEEEauthorblockN{First A. Author}
\IEEEauthorblockA{Department of Electrical Engineering\\
University Name\\
City, Country\\
email@university.edu}
\and
\IEEEauthorblockN{Second B. Author}
\IEEEauthorblockA{Research Laboratory\\
Company / Institute\\
City, Country\\
email@company.com}
}

\maketitle

\begin{abstract}
This skeleton demonstrates the structure of an IEEE conference paper using the \texttt{IEEEtran} class. Replace the sample text with your manuscript content. Keep the abstract within the page limit required by the target venue.
\end{abstract}

\begin{IEEEkeywords}
LaTeX, IEEEtran, conference paper, template
\end{IEEEkeywords}

\IEEEpeerreviewmaketitle

\section{Introduction}
IEEE papers typically begin with motivation, related work summary, and a clear statement of contributions.

\section{Method}
Describe the proposed method with equations, algorithms, and figures.

\begin{equation}
    E = mc^2
    \label{eq:sample}
\end{equation}

Equation~\eqref{eq:sample} is only a placeholder.

\section{Results}
Present experimental settings, baselines, and quantitative results.

\begin{table}[!t]
\caption{Sample Results}
\label{tab:sample}
\begin{center}
\begin{tabular}{lcc}
\hline
Method & Accuracy & Latency (ms) \\
\hline
Baseline & 90.1 & 12.4 \\
Ours & 93.5 & 10.8 \\
\hline
\end{tabular}
\end{center}
\end{table}

\section{Discussion}
Summarize implications, limitations, and future work without over-claiming.

\section{Conclusion}
Restate the problem, the approach, and the main experimental finding in a few sentences.

\section*{Acknowledgment}
This work was supported in part by Grant No.~XXXX.

\begin{thebibliography}{00}
\bibitem{ref1}
A.~Author, ``Title of paper,'' \textit{IEEE Trans.~XXX}, vol.~XX, no.~XX, pp.~XXX--XXX, Year.

\bibitem{ref2}
B.~Author, ``Title of conference paper,'' in \textit{Proc. IEEE Conf.}, City, Year, pp.~XXX--XXX.
\end{thebibliography}

\end{document}
